\chapter{Pell's Equation}

\section{Approximations of \texorpdfstring{$\sqrt{2}$}{√2}}

\begin{enumerate}
\item Discuss Pythagorean triples of the form $(a,a+1,c)$.
\item Connect these Pythagorean triples to the equation $x^2-2y^2=-1$.
\item Use the $(3,4,5)$ triangle to find the first solution.
\item Use the first solution to generate all solutions.
\item List the pythagoran triples of the form $(a,a+1,c)$ and their resulting approximations to $\sqrt{2}$.
\item Observe that it is possible to approach this subject with continued fractions.
\end{enumerate}

% \section{Algebraic Fields}

% \begin{definition}
%   An \emph{algebraic number} is a solution $\xi$ to a polynomial equation
%   \begin{equation*}
%     a_n\xi^n+\cdots+a_1\xi+a_0=0
%   \end{equation*}
%   with integer coefficients. An algebraic number $\xi$ is said to be of \emph{degree $n$} if $\xi$ is the root of a polynomial of degree $k$ if and only if $n\leq k$.
% \end{definition}

% \begin{definition}
%   Consider an algebraic number $\xi$. The \emph{algebraic field} $\mathbb{Q}(\xi)$ consists of all numbers of the form
%   \begin{equation*}
%     \frac{P(\xi)}{Q(\xi)}
%   \end{equation*}
%   where $P$ and $Q$ are polynomials with integer coefficients, and $Q(\xi)$ is nonzero.
% \end{definition}

\section{Quadratic Fields}

In this section, we assume that $m$ is a square-free integer.

\begin{definition}
  We define $\mathbb{Z}[\sqrt{m}]$ to be the set of numbers of the form
  \begin{equation*}
    a+b\sqrt{m}.
  \end{equation*}
  Addition and multiplication are defined on $\mathbb{Z}[\sqrt{m}]$ by
  \begin{align*}
    (a+b\sqrt{m})+(c+d\sqrt{m}) & := (a+c)+(b+d)\sqrt{m}, \\
    (a+b\sqrt{m})(c+d\sqrt{m}) & := (ac + mbd) + (ad + bc)\sqrt{m}.
  \end{align*}
\end{definition}

\begin{definition}
  Consider a square-free integer $m$. We define the \emph{quadratic field} $\mathbb{Q}(\sqrt{m})$ to consist of numbers of the fractions in $\mathbb{Z}[\sqrt{m}]$; that is, it consists of numbers of the form
  \begin{equation*}
    \frac{a+b\sqrt{m}}{c+d\sqrt{m}},
  \end{equation*}
  where $c+d\sqrt{m}\neq 0$. Numbers in $\mathbb{Q}(\sqrt{m})$ can be added and multiplied by the following rules:
  \begin{align*}
    \frac{a+b\sqrt{m}}{c+d\sqrt{m}}+\frac{e+f\sqrt{m}}{g+h\sqrt{m}} & := \frac{(a+b\sqrt{m})(g+h\sqrt{m})+(e+f\sqrt{m})(c+d\sqrt{m})}{(c+d\sqrt{m})(g+h\sqrt{m})}, \\
    \frac{a+b\sqrt{m}}{c+d\sqrt{m}} \frac{e+f\sqrt{m}}{g+h\sqrt{m}} & := \frac{(a+b\sqrt{m})(e+f\sqrt{m})}{(c+d\sqrt{m})(g+h\sqrt{m})}
  \end{align*}
\end{definition}

\section{Solutions to Pell's Equation}

%%% Local Variables:
%%% mode: latex
%%% TeX-master: "elementary-number-theory"
%%% TeX-engine: luatex
%%% End:
